\section{Sources et bibliographie}
\label{annexe-bibliographie}
\begingroup
\footnotesize
\urlstyle{same}
\newenvironment{biblist}{%
  \begin{list}{}{%
    \setlength{\leftmargin}{1.5em}%
    \setlength{\itemindent}{-1.5em}%
    \setlength{\labelwidth}{0pt}%
    \setlength{\labelsep}{0pt}%
    \setlength{\itemsep}{0.5pt plus 0.5pt}%
    \setlength{\parsep}{0pt}%
    \setlength{\topsep}{2pt}%
    \setlength{\partopsep}{0pt}}%
  \raggedright}{\end{list}}

\noindent La présente bibliographie réunit toutes les références citées dans le commentaire historique, les introductions, les annexes et les dossiers de recherche de l'édition (dossiers R1 à R8, conservés avec les fichiers de travail). Les ouvrages repérés mais non consultés, ou cités de seconde main, sont signalés comme tels. Les adresses électroniques ont été consultées les 17 et 18~octobre 2026, sauf indication contraire. Les cotes du «~dépôt~» renvoient aux dossiers numériques des archives familiales (\texttt{archives\_militaires/} et \texttt{biographie\_et\_famille/}).

%======================================================================
\subsection*{Abréviations bibliographiques}

\begin{multicols}{2}
\begin{biblist}
\item \textbf{AD} : Archives départementales.
\item \textbf{AD Somme, Espel 2018} : \textsc{Espel} 2018 (III).
\item \textbf{AIF Project} : UNSW Canberra, \textit{The AIF Project} (IV).
\item \textbf{Aitken 1916} : \textsc{Aitken}, \textit{Canada in Flanders} (II\,A).
\item \textbf{Baedeker 1912} : \textsc{Baedeker}, \textit{Le Sud-Ouest de la France} (II\,B).
\item \textbf{Baedeker 1913} : \textsc{Baedeker}, \textit{Paris and Environs} (II\,B).
\item \textbf{Baedeker 1914} : \textsc{Baedeker}, \textit{Southern France} (II\,B).
\item \textbf{Banks et Chell 1924}, \textbf{Banks \& Chell 1924} : \textit{With the 10th Essex in France} (II\,B).
\item \textbf{Barbé 1917} : \textsc{Barbé}, «~L'hiver 1916-1917~» (II\,D).
\item \textbf{Bean 1941} : \textsc{Bean}, \textit{The AIF in France, 1916} (II\,A).
\item \textbf{Becke 1938} (une fois, par erreur, \textbf{Becke 1945}, date de la part~3B) : \textsc{Becke}, \textit{Order of Battle of Divisions}, part 3A (II\,A).
\item \textbf{Beds War Diary} : voir \textbf{JMO 7/Bedf}.
\item \textbf{Bewsher 1921} : \textit{The History of the 51st (Highland) Division} (II\,B).
\item \textbf{Blackwood's}, juin 1917 : «~Claims~», \textit{Blackwood's Magazine} (II\,D).
\item \textbf{Blanchard 1906} : \textsc{Blanchard}, \textit{La Flandre} (II\,B).
\item \textbf{Bulletin des armées} : \textit{Bulletin des armées de la République} (II\,D).
\item \textbf{Burnichon 1916} : \textit{La Compagnie de Jésus en France} (II\,B).
\item \textbf{Cassini/EHESS} : \textit{Des villages de Cassini aux communes d'aujourd'hui} (IV).
\item \textbf{CFD Somme} : Wikipédia, «~Chemins de fer départementaux de la Somme~» (IV).
\item \textbf{Clifton et Grimaux 1883} : \textit{Nouveau dictionnaire anglais-français} (II\,C).
\item \textbf{Combes 2020}, \textbf{PC-2020} : P.~\textsc{Combes}, étude familiale illustrée, 2020--2021 (I\,A).
\item \textbf{Cronier 2013} : \textit{Permissionnaires dans la Grande Guerre} (III).
\item \textbf{CWGC} : Commonwealth War Graves Commission (IV).
\item \textbf{Daraux 2002} : \textsc{Daraux}, «~L'œuvre de Louis Mazetier~» (III).
\item \textbf{Delany 1907} : \textit{The Catholic Encyclopedia}, «~St. Benedict Joseph Labre~» (II\,B).
\item \textbf{Dictionary 1919}, \textbf{French-English Dictionary 1919} : \textit{French-English Dictionary of Military Terms} (II\,C).
\item \textbf{Ellis [1920]} : \textit{The Story of the Fifth Australian Division} (II\,B).
\item \textbf{Faidherbe 1871} : \textit{Campagne de l'armée du Nord} (II\,B).
\item \textbf{Falls 1933} : \textsc{Falls}, \textit{Military Operations, Macedonia}, vol.~I (II\,A).
\item \textbf{Feller-Vest 2015} : \textit{DHS}, «~Commynes, Philippe de~» (III).
\item \textbf{Fernand-Laurent 1917} : \textit{Chez nos alliés britanniques} (II\,B) ; \textbf{Tessan, dans Fernand-Laurent 1917} : sa préface.
\item \textbf{Fortescue 1928} : \textit{A Short Account of Canteens} (II\,B).
\item \textbf{Fraser et Gibbons 1925} : \textit{Soldier and Sailor Words and Phrases} (II\,C).
\item \textbf{FSPB 1914} : \textit{Field Service Pocket Book, 1914} (II\,A).
\item \textbf{Haig 1919} : \textsc{Boraston} (éd.), \textit{Sir Douglas Haig's Despatches} (II\,A).
\item \textbf{Henson, journal} : Durham University, \textit{Henson Journals} (I\,B).
\item \textbf{IGN} : Institut national de l'information géographique, Géoplateforme (IV).
\item \textbf{INSEE} : fichier des personnes décédées ; répertoire SIRENE (IV).
\item \textbf{IWM Q 961} : Imperial War Museum, photographie Q~961 (IV).
\item \textbf{JMO 7/Bedf} (introductions : \textbf{JMO du 7th Bedfordshire}) : \textit{War Diary} du 7th (Service) Battalion, Bedfordshire Regiment, Bedfordshire Archives Service, X550/8/1 (I\,B).
\item \textbf{Jones 1928}, \textbf{Jones 1931} : \textsc{Jones}, \textit{The War in the Air}, vol.~II et III (II\,A).
\item \textbf{Kipling 1923} : \textit{The Irish Guards in the Great War} (II\,B).
\item \textbf{Lafon 1922} : \textit{Les Chemins de fer français pendant la guerre} (II\,B).
\item \textbf{LLT} : C.~\textsc{Baker}, \textit{The Long, Long Trail} (IV).
\item \textbf{London Gazette} : \textit{The London Gazette} et ses suppléments (II\,D).
\item \textbf{March 1925} : \textit{Mouvement des prix et des salaires pendant la guerre} (II\,B).
\item \textbf{Maurois 1920} : \textit{The Silence of Colonel Bramble} (II\,B).
\item \textbf{Mérimée}, \textbf{Palissy} : bases du ministère de la Culture (IV).
\item \textbf{Michelin 1919} : \textit{The Somme}, vol.~1 ; \textbf{Michelin 1919, liste des hôtels} : \textit{Amiens before and during the War} ; \textbf{Michelin 1921} : \textit{St-Quentin -- Cambrai} (II\,B).
\item \textbf{Miles 1938} : voir \textbf{OH 1916-II}.
\item \textbf{Nichols 1922} : \textit{The 18th Division in the Great War} (II\,B).
\item \textbf{Nicholson 1962} : \textit{Canadian Expeditionary Force} (II\,A).
\item \textbf{Nicot 1972} : \textsc{Nicot}, \textsc{Schillinger}, \textsc{Ficat}, \textit{Inventaire sommaire} (III).
\item \textbf{Offenstadt 2004} : \textit{Le Chemin des Dames} (III).
\item \textbf{OH} : \textit{Official History}, histoire officielle britannique (\textit{History of the Great War based on Official Documents}).
\item \textbf{OH 1916-I} (introductions : \textit{OH 1916}, I) : \textsc{Edmonds} 1932 (II\,A).
\item \textbf{OH 1916-II} (introductions : \textit{OH 1916}, II) : \textsc{Miles} 1938 (II\,A).
\item \textbf{OH 1917-I} (introductions : \textit{OH 1917}, I) : \textsc{Falls} 1940 (II\,A).
\item \textbf{OQLF, GDT} : Office québécois de la langue française, \textit{Grand Dictionnaire terminologique} (IV).
\item \textbf{OSM} : OpenStreetMap, service Nominatim (IV).
\item \textbf{Partridge 1937} : \textit{A Dictionary of Slang} (II\,C).
\item \textbf{PC-Annexe} : P.~\textsc{Combes}, livre familial, \texttt{Combes-annexe} (annexe~2, «~Les Calmels~») (I\,A).
\item \textbf{PC-Bio} : P.~\textsc{Combes}, «~Mon grand-père Hippolyte Combes (et ses sœurs)~» (I\,A) ; introductions : P.~Combes, «~Mon grand-père~».
\item \textbf{PC-Livre} : P.~\textsc{Combes}, \textit{Une étude de la famille Combes de Castelnau-Pégayrols} (I\,A) ; introductions : P.~Combes, \textit{Étude}.
\item \textbf{Pouvillon 1938} : discours de réception, \textit{Recueil de l'Académie de Montauban} (II\,D).
\item \textbf{Puckle 1917} : \textit{The Army Service Corps of the British Army} (II\,A).
\item \textbf{R1}\dots\textbf{R8}, \textbf{F}\textit{n} : dossiers de recherche de l'édition et numéro du fait (p.~ex. «~R5, F61~»).
\item \textbf{RATP} : Régie autonome des transports parisiens (IV).
\item \textbf{Sassoon 1930} : \textit{Memoirs of an Infantry Officer} (II\,B).
\item \textbf{Shakespear 1921} : \textit{The Thirty-Fourth Division} (II\,B).
\item \textbf{SHD} : Service historique de la Défense, Vincennes (I\,B).
\item \textbf{Statistics 1922} : War Office, \textit{Statistics of the Military Effort} (II\,A).
\item \textbf{Tattersfield 2021} : «~The workhorse of the staff~» (III).
\item \textbf{Thompson 1990} : \textit{The Free Church Army Chaplain} (III).
\item \textbf{TNA} : The National Archives, Kew (I\,B).
\item \textbf{UNESCO 943-033} : Centre du patrimoine mondial, «~Beffrois de Belgique et de France~» (IV).
\item \textbf{Veterinary Services 1925} : \textsc{Blenkinsop} et \textsc{Rainey} (éd.) (II\,A).
\item \textbf{Ville de Blois 2015} : dossier de presse «~Les belles images de Poulain~» (IV).
\item \textbf{Ward 2022} : «~Wren's Shoe Polish and Northampton~» (III).
\item \textbf{Whittem et Long 1917} : \textit{French for Soldiers} (II\,B).
\item \textbf{Wyrall [1926]} : \textit{The Die-Hards in the Great War} (II\,B).
\item \textbf{Yapp 1918} : \textit{The Romance of the Red Triangle} (II\,B).
\item \textbf{Zeller 1946} : «~Procès à reviser ?~» (III).
\item \textbf{Zyromski 1921} : \textit{Maurice de Guérin} (II\,B).
\end{biblist}
\end{multicols}

%======================================================================
\subsection*{I. Sources manuscrites et archives}

\subsubsection*{A. Archives familiales (dépôt numérique de la famille Combes)}

\begin{biblist}
\item \textbf{Le \textit{Livre-Journal}.} \textsc{Combes}, Hippolyte, \textit{Livre-Journal}, carnet manuscrit, 25~juillet 1916 -- 30~avril 1917, 93~feuillets, collection familiale ; 186~photographies du 26~novembre 2025 (\texttt{scans\_haute\_qualite/}).
\item \texttt{archives\_militaires/fiche\_et\_parcours/}. Fiche matricule d'Hippolyte Combes, recrutement de Toulouse, classe 1896, n\textsuperscript{o}~911, établie le 30~novembre 1921, copie numérique (\texttt{Fiche\_militaire\_Combes\_Hippolyte.jpg}) ; courriels de Pierre Combes (11~octobre 2021 ; 2024 ; «~au sujet d'Hippolyte, parcours militaire~») ; échange WhatsApp avec Pierre Combes (26~novembre 2025 ; 24~mai 2026) ; D.~\textsc{Boyer}, «~H. Combes interprète du prince de Galles~», courriel du 7~décembre 2025.
\item \texttt{archives\_militaires/lettres\_de\_guerre/}. Lettre d'Hippolyte Combes à un capitaine, «~ancien camarade à la Faculté de Toulouse~», «~Au front, le 25~février 1916~», copie autographe, 3~p. (\texttt{lettre fev 1916 au front à un capitaine.pdf} ; annexe, doc.~1) ; lettre à sa sœur, «~Au front, vendredi 16~juin 1916~», 4~p. (\texttt{lettre de Hippolyte à sa soeur Josephine - 16 juin 1916 - Au front.pdf}) ; minutes du rapport sur l'incident de Morlancourt («~L'interprète Combes au sous-lieutenant\dots~») et de la supplique sur l'indemnité des interprètes, carnet quadrillé, p.~26--27 et 37 (\texttt{lettre de l'interprete combes pour un incident.pdf}) ; poème anglais (\texttt{un poeme.pdf} ; annexe, doc.~4) ; carte-lettre à Marguerite, Paris, 25~novembre 1918 (\texttt{lettre à bonne maman 25 novembre 1918.jpg}) ; cartes à Marguerite, Paris, 26~novembre 1918, et Les Andelys, 28~novembre et 2~décembre 1918 ; cinq clichés complémentaires (\texttt{20250413\_*.jpg}) ; projet d'ouvrage autographe \textit{Feuillets de Route — Souvenirs d'un Interprète auprès de l'armée britannique} (feuillet de titre, feuillet de division, sommaire en vingt-deux chapitres et deux états du chapitre premier, 5~feuillets, 6~pages écrites, \texttt{IMG\_20261022\_0001.pdf}, communiqué par Pierre Combes le 22~octobre 2026 ; annexe, doc.~6).
\item \texttt{archives\_militaires/photos\_militaires/}. Photographies de l'école d'artillerie de Toulon, avril 1915 (dont le cliché de Marius Bar, «~Toulon --- 10\textsuperscript{e} d'Artillerie. - 39\textsuperscript{e} Batterie~») ; «~collègue Nogent le Rotrou~» ; «~Angèle Brillouet~» ; école militaire des Andelys ; camarades, Ailly-sur-Somme, 9~juin 1919 ; portraits en uniforme.
\item \texttt{archives\_militaires/1st\_Irish\_Guards/} et \texttt{Army\_Service\_Corps/}. Copie du journal de marche TNA, WO~95/1216/1 (quatre fichiers PDF) et notes sur R.~Kipling, \textit{The Irish Guards in the Great War} ; \textit{The Ringing World}, n\textsuperscript{o}~212, 9~avril 1915 (PDF).
\item \texttt{biographie\_et\_famille/etat\_civil\_et\_actes/}. Acte de naissance de Marie Lucien Hippolyte Combes, Millau, n\textsuperscript{o}~395, 16~août 1876 (\texttt{1876\_acte\_naissance\_hippolyte\_Combes.png}) ; acte de naissance de Marie Marguerite Calmels, Millau, n\textsuperscript{o}~137, 16~avril 1892 (\texttt{1892\_acte\_naissance\_Marguerite\_Calmels.png}) ; faire-part de décès de M\textsuperscript{me} Jeanne-Eulalie Monestier, veuve Cartairade (Montpellier, 2~avril 1888) et de M\textsuperscript{me} Guillaume Cartairade, née Émilie Noel (Lisside, Gard, 4~mars 1889 ; Millau, imp. Artières et J.~Maury) ; menu manuscrit du repas de noces Combes--Calmels, «~Menu --- 25~octobre 1917~», et photographie du mariage (Millau, 1917) ; mariage de Valentine Combes et d'André Combès, Toulouse, 7~septembre 1946.
\item \texttt{biographie\_et\_famille/biographie/}. \textsc{Combes}, Pierre, «~Mon grand-père Hippolyte Combes (et ses sœurs)~», biographie familiale dactylographiée, inédite, s.~d. (\texttt{.docx} et \texttt{.txt}) [PC-Bio] ; \textsc{Combes}, Pierre, étude familiale illustrée (arbres généalogiques), version 2020--2021, scan de 64~vues (\texttt{Pierre\_Combes\_famille\_Combes\_2020.pdf} : arbres p.~52, 96, 98, 106, 107, 114) [PC-2020] ; du même, \texttt{Pierre\_Combes\_famille\_Bonnel\_Goubault.pdf} et \texttt{fonthugne de rayrols.pdf} ; notes \texttt{Calmels.docx}, \texttt{photosPierreCombes.docx} (légendes), \texttt{Hipollyte 1 et Francois.docx}, \texttt{Madeleine Calmels.pdf}.
\item \texttt{biographie\_et\_famille/Livre\_Pierre\_Combes/}. \textsc{Combes}, Pierre, \textit{Une étude de la famille Combes de Castelnau-Pégayrols}, dactylographie inédite, version du 19~mars 2023, fichiers \texttt{Combes-1} à \texttt{Combes-4} [PC-Livre] et \texttt{Combes-annexe} (annexe~2, «~Les Calmels~») [PC-Annexe].
\item \texttt{biographie\_et\_famille/correspondance\_familiale/}. Lettre de Françoise Combes, Périgueux, 26~janvier 191[3] ; lettres d'Hippolyte Combes, Toulouse, 17~septembre 1949 et 23~octobre 1961 (papier à en-tête du Caousou) ; dédicace de tante Françoise, juillet 1936 ; lettre («~lettre-Maman~»).
\item \texttt{biographie\_et\_famille/discours\_et\_academie\_montauban/}. \textit{Lycée Ingres. Distribution des prix du 13~juillet 1928. Discours prononcé par M.~Hippolyte Combes, professeur d'anglais}, brochure dédicacée à son fils Jean (5~vues) ; \textit{Recueil de l'Académie de Montauban}, 1932 et 1938 (voir II\,D) ; \textit{Mémoires de la Société des lettres} [\dots] \textit{de l'Aveyron}, 1865 (H.~Combes, avocat démissionnaire).
\item \texttt{biographie\_et\_famille/La\_Valette/} et \texttt{portraits\_et\_lieux/}. Oncle François, histoire manuscrite de La Valette (23~vues) et note «~La Valette, Pierre et François Combes~» ; portraits et vues (Hippolyte Combes à 21~ans, 1897 ; au Caousou ; à Cambridge, années 1930 ; Marguerite Calmels ; maison Combes de Millau ; 55, faubourg du Moustier, Montauban ; Carmel de Périgueux ; fratrie Combes).
\item \textbf{Autres pièces citées.} Minute de lettre à un «~Monsieur le Président~» non identifié, s.~d., carnet quadrillé, p.~37 ; courriels de P.-A.~Hauseux (23~novembre 2025 ; 18 et 22~octobre 2026), de Pierre Combes (26~novembre 2025 ; 18 et 22~octobre 2026) et de Jacques Vincens (18~octobre 2026).
\end{biblist}

\subsubsection*{B. Archives publiques}

\begin{biblist}
\item \textbf{Archives départementales de l'Aveyron.} État civil de Millau : actes de naissance de 1876 et 1892 (copies au dépôt, voir I\,A) ; acte de mariage du 25~octobre 1917 et transcription du décès de Joséphine Combes (1914) : non consultés.
\item \textbf{Archives départementales d'Eure-et-Loir.} État civil de Nogent-le-Rotrou, décès de 1914 : non consulté (accès en ligne bloqué, \url{https://archives28.fr/}).
\item \textbf{Archives départementales de la Haute-Garonne.} Registre matricule du recrutement de Toulouse, classe 1896, n\textsuperscript{o}~911 (Combes, Marie Lucien Hippolyte) ; copie au dépôt.
\item \textbf{Archives départementales de la Somme.} Fonds de l'Office notarial de Toutencourt, aujourd'hui Acheux-en-Amiénois (1729--1904) : voir \textsc{Espel} 2018 (III).
\item \textbf{Archives du Carmel de Lisieux.} «~Supplique des soldats envoyée au pape pour la béatification de Thérèse~», \url{https://archives.carmeldelisieux.fr}.
\item \textbf{Bedfordshire Archives Service} (Bedford), X550/8/1 : \textit{War Diary}, 7th (Service) Battalion, Bedfordshire Regiment (54th Infantry Brigade, 18th Division) [JMO 7/Bedf], transcriptions intégrées au catalogue en ligne (\url{https://bedsarchivescat.bedford.gov.uk/Details/archive/}, suivi du numéro de notice) : juillet 1916, /1607 (110139701) et /1607b (110153196) ; août 1916, /1608 (110153222), avec l'ordre d'opérations n\textsuperscript{o}~26 du 23~août ; septembre 1916, /1609 (110139707), avec les ordres n\textsuperscript{os}~28 à 31 ; octobre 1916, /1610 (110139708) ; décembre 1916, /1612 (110153355) ; janvier 1917, /1701 (110153356) ; février 1917, /1702 (110139700) ; mars 1917, /1703 (110153357) ; avril 1917, /1704 (110153358), ordres n\textsuperscript{os}~4 à 6 ; notice du bataillon (110003793). Novembre 1916 (/1611) non retrouvé en ligne.
\item \textbf{Durham University.} \textsc{Henson}, Herbert Hensley, \textit{Journals}, vol.~20, p.~424, 28~août 1916, \url{https://henson.durham.ac.uk/journals/16-08-28}.
\item \textbf{Service historique de la Défense} (Vincennes), non consultés : GR~16~N~509 (GQG, interprètes : affectations, mutations) ; GR~17~N (Mission militaire française attachée à l'armée britannique, cartons 295--486), notamment 17~N~295 (organisation), 360--361 (5\textsuperscript{e} armée britannique), 393--396 (population civile, réquisitions, réclamations, achats), 398--399 (unités françaises, dont le 19\textsuperscript{e} escadron du train), 405--406 (soldes et indemnités), 446 (agriculture, battages), 447--448 (listes nominatives), 463--479 (fiches individuelles et carnets de notes des interprètes).
\item \textbf{The National Archives} (Kew). \textit{War Diaries} : WO~95/1216/1 (copie au dépôt, voir I\,A) ; non consultés : WO~95/2032/2 (18th Division, Divisional Train, 150--153 Companies ASC, juillet 1915 -- avril 1919, \url{https://discovery.nationalarchives.gov.uk/details/r/C7352984}), WO~95/2015/2, 2015/3 et 2016/1 (18th Division, General Staff), WO~95/2018/2 (Adjutant and Quarter-Master General), WO~95/2041/2 à 2041/5 (QG de la 54th Infantry Brigade). Dossiers d'officiers (notices du catalogue Discovery) : WO~339/16884 (W.~H.~Levy) ; WO~339/21191 (Captain Edward Maurice West, Royal Army Service Corps, 1914--1919) ; WO~339/25186 (H.~F.~Hughes) ; WO~339/28396 (E.~M.~Ford) ; WO~339/44046 (Lieutenant G.~E.~Bradfield, RASC) ; WO~374/41727 (R.~Leonard) ; WO~374/41729 (Captain Robert Marsh Leonard, ASC) ; requêtes dans les séries WO~339 et WO~374 («~Curtis~», «~Gilbert~», «~Panton~», «~Williams~», avec «~Army Service Corps~»).
\end{biblist}

%======================================================================
\subsection*{II. Sources imprimées}

\subsubsection*{A. Histoires officielles, publications et textes officiels}

\begin{biblist}
\item \textsc{Aitken}, Sir Max [Lord Beaverbrook], \textit{Canada in Flanders. The Official Story of the Canadian Expeditionary Force}, vol.~I, préface d'A.~Bonar Law, introduction de Sir Robert Borden, Londres, Toronto et New York, Hodder \& Stoughton, 1916 ; vol.~II, Londres, Hodder \& Stoughton, 1917 (préface du 27~janvier 1917). \url{https://archive.org/details/canadainflanders01beavuoft} ; \url{https://archive.org/details/officialstory00beavuoft} ; Project Gutenberg n\textsuperscript{o}~46114.
\item \textsc{Bean}, C.~E.~W., \textit{The Australian Imperial Force in France, 1916} (\textit{Official History of Australia in the War of 1914--1918}, vol.~III), 12\textsuperscript{e}~éd., Sydney, Angus \& Robertson, 1941 (chap.~XIII, AWM RCDIG1069465, \url{https://www.awm.gov.au/collection/RCDIG1069465}).
\item \textsc{Becke}, A.~F., \textit{Order of Battle of Divisions}, part~3A : \textit{New Army Divisions (9--26)} (\textit{History of the Great War based on Official Documents}), Londres, HMSO, 1938. \url{https://archive.org/details/order-battle-divisions-p-3-a}.
\item \textsc{Blenkinsop}, L.~J., et \textsc{Rainey}, J.~W. (éd.), \textit{Veterinary Services} (\textit{History of the Great War based on Official Documents}), Londres, HMSO, 1925. \url{https://archive.org/details/veterinary-services-1925}.
\item \textsc{Boraston}, J.~H. (éd.), \textit{Sir Douglas Haig's Despatches (December 1915 -- April 1919)}, Londres et Toronto, J.~M.~Dent \& Sons, 1919 (dépêches du 23~décembre 1916 et du 31~mai 1917). \url{https://archive.org/details/sirdouglashaigsd00haiguoft}.
\item \textit{Code de commerce}, Paris, 1807, livre~I, titre~II, «~Des livres de commerce~», art.~8.
\item \textsc{Edmonds}, J.~E. (comp.), \textit{Military Operations, France and Belgium, 1916}, vol.~I : \textit{Sir Douglas Haig's Command to the 1st July : Battle of the Somme}, cartes d'A.~F.~Becke, Londres, Macmillan, 1932 [OH 1916-I]. \url{https://archive.org/details/in.ernet.dli.2015.210679}.
\item \textsc{Falls}, Cyril (comp.), \textit{Military Operations, Macedonia}, vol.~I : \textit{From the Outbreak of War to the Spring of 1917}, Londres, HMSO, 1933. \url{https://archive.org/details/military-operations-macedonia-vol-1}.
\item \textsc{Falls}, Cyril (comp.), \textit{Military Operations, France and Belgium, 1917}, vol.~I : \textit{The German Retreat to the Hindenburg Line and the Battles of Arras}, Londres, Macmillan, 1940 [OH 1917-I]. \url{https://archive.org/details/in.ernet.dli.2015.210683}.
\item \textit{Field Service Pocket Book, 1914}, General Staff, War Office, Londres, HMSO, 1914 [FSPB 1914]. \url{https://archive.org/details/b28998558}.
\item \textit{Handbook to the Most Excellent Order of the British Empire}, éd. A.~Winton Thorpe, Londres, 1921. \url{https://archive.org/details/handbooktomostex0000awin}.
\item Indult du cardinal Caprara pour la réduction des fêtes, 9~avril 1802, publié par l'arrêté des consuls du 29~germinal an~X (19~avril 1802).
\item \textsc{Jones}, H.~A., \textit{The War in the Air, being the Story of the Part played in the Great War by the Royal Air Force}, vol.~II, Oxford, Clarendon Press, 1928 ; vol.~III, 1931 (app.~XII). \url{https://archive.org/details/warinairbeingsto02rale} ; \url{https://archive.org/details/warinairbeingsto03rale}.
\item Loi du 3~juillet 1877 relative aux réquisitions militaires.
\item \textsc{Miles}, Wilfrid (comp.), \textit{Military Operations, France and Belgium, 1916}, vol.~II : \textit{2nd July 1916 to the End of the Battles of the Somme}, préface de J.~E.~Edmonds, cartes d'A.~F.~Becke, Londres, Macmillan, 1938, et volume \textit{Maps and Appendices}, 1938 [OH 1916-II]. \url{https://archive.org/details/in.ernet.dli.2015.210680} ; \url{https://archive.org/details/in.ernet.dli.2015.210678}.
\item \textsc{Nicholson}, G.~W.~L., \textit{Canadian Expeditionary Force 1914--1919}, Ottawa, Queen's Printer, 1962.
\item \textsc{Puckle}, F.~K., \textit{The Army Service Corps of the British Army and the Organization of the Transport and Transportation at the Front in France. Lectures delivered before the officers of the Quartermaster Corps and Quartermaster Reserve Corps, at Washington, D.C., on May 2 and May 9, 1917}, Washington, Government Printing Office, 1917. \url{https://archive.org/details/armyservicecorps00puck}.
\item \textit{Report of the Ministry, Overseas Military Forces of Canada, 1918}, Londres, 1919. \url{https://archive.org/details/reportofministr00cana}.
\item \textsc{War Office}, \textit{Statistics of the Military Effort of the British Empire during the Great War, 1914--1920}, Londres, HMSO, 1922. \url{https://archive.org/details/statisticsofmili00grea}.
\end{biblist}

\subsubsection*{B. Histoires d'unités, témoignages, guides et ouvrages anciens}

\begin{biblist}
\item \textit{19\textsuperscript{e} Escadron du train des équipages militaires. Historique. Guerre 1914--1918}, Paris, Lavauzelle, 1921 (non consulté).
\item \textsc{Atkinson}, C.~T., \textit{The Queen's Own Royal West Kent Regiment, 1914--1919}, Londres, Simpkin, Marshall, Hamilton, Kent \& Co., 1924. \url{https://archive.org/details/dli.ministry.19498}.
\item \textsc{Baedeker}, Karl, \textit{Le Sud-Ouest de la France, de la Loire à la frontière d'Espagne. Manuel du voyageur}, Leipzig, K.~Baedeker, 1912. \url{https://archive.org/details/lesudouestdelafr00karl}.
\item \textsc{Baedeker}, Karl, \textit{Paris and Environs, with Routes from London to Paris}, 18\textsuperscript{e}~éd., Leipzig, K.~Baedeker, 1913. \url{https://archive.org/details/parisenviro00karl}.
\item \textsc{Baedeker}, Karl, \textit{Southern France including Corsica. Handbook for Travellers}, 6\textsuperscript{e}~éd., Leipzig, K.~Baedeker, 1914. \url{https://archive.org/details/southernfrancein00karl}.
\item \textsc{Banks}, T.~M., et \textsc{Chell}, R.~A., \textit{With the 10th Essex in France}, Londres, Gay \& Hancock, 1\textsuperscript{re}~éd. 1921 ; 2\textsuperscript{e}~éd. 1924 (pagination citée). \url{https://archive.org/details/10thessexfrance}.
\item \textsc{Barthès}, Émile, \textit{Eugénie de Guérin d'après des documents inédits}, t.~II : \textit{Après la mort de son frère Maurice}, thèse de doctorat ès lettres, Paris, Lecoffre-J.~Gabalda, et Albi, Imprimerie coopérative du Sud-Ouest, 1929. \url{https://archive.org/details/IA41551824_0003}.
\item \textsc{Bewsher}, F.~W., \textit{The History of the 51st (Highland) Division, 1914--1918}, Édimbourg et Londres, W.~Blackwood \& Sons, 1921. \url{https://archive.org/details/historyof51sthig00bews}.
\item \textsc{Blanchard}, Raoul, \textit{La Flandre. Étude géographique de la plaine flamande en France, Belgique et Hollande}, Paris, Armand Colin, 1906. \url{https://archive.org/details/laflandretudegog00blan}.
\item \textsc{Burnichon}, Joseph, \textit{La Compagnie de Jésus en France. Histoire d'un siècle, 1814--1914}, t.~II (1830--1845), Paris, G.~Beauchesne, 1916 ; t.~III (1845--1860), 1919. \url{https://archive.org/details/lacompagniedej02burn} ; \url{https://archive.org/details/lacompagniedej03burn}.
\item \textsc{Calvet}, Arthur, \textit{Le Père Paul Ginhac de la Compagnie de Jésus}, 5\textsuperscript{e}~éd., Tournai, Casterman, et Paris, Ch.~Amat, 1916. \url{https://archive.org/details/leperepaulginhac0000lepe}.
\item \textsc{Chesterman}, Hugh, «~With a Train Company in France, 1914--1918~», \textit{The Nineteenth Century and After}, 86, n\textsuperscript{o}~511, septembre 1919, p.~571.
\item \textsc{Cortyl}, E., «~La dentelle à Bailleul~», \textit{Bulletin du Comité flamand de France}, 1903, p.~225--237 (cité d'après Blanchard ; non consulté).
\item \textsc{Delany}, Joseph F., «~St. Benedict Joseph Labre~», \textit{The Catholic Encyclopedia}, vol.~2, New York, Robert Appleton Company, 1907. \url{https://www.newadvent.org/cathen/02442a.htm}.
\item \textsc{Ellis}, A.~D., \textit{The Story of the Fifth Australian Division, being an Authoritative Account of the Division's Doings in Egypt, France and Belgium}, Londres, Hodder \& Stoughton, [1920]. \url{https://archive.org/details/storyoffifthaust00ellirich}.
\item \textsc{Faidherbe}, Louis, \textit{Campagne de l'armée du Nord en 1870-1871}, Paris, E.~Dentu, 1871.
\item \textsc{Fernand-Laurent}, \textit{Chez nos alliés britanniques (With our British Allies in the Field). Notes et souvenirs d'un interprète}, préface de François de Tessan, Paris, Boivin, 1917. \url{https://archive.org/details/cheznosallisbr00fernuoft}.
\item \textsc{Fleurance}, Gustave de, \textit{Expulseurs et expulsés}, préface d'É.~Drumont, Paris, Letouzey et Ané, [1888]. \url{https://archive.org/details/expulseursetexpu00fleu}.
\item \textsc{Fortescue}, Sir John W., \textit{A Short Account of Canteens in the British Army}, Cambridge, Cambridge University Press, 1928. \url{https://archive.org/details/canteensbritisharmy}.
\item \textsc{Greenhalgh}, Elizabeth (éd.), \textit{Liaison : General Pierre des Vallières at British General Headquarters, January 1916 to May 1917}, Army Records Society, vol.~36, 2016.
\item \textit{Hautes payes de guerre, indemnités de combat, primes d'alimentation, pécules}, Paris, L.~Fournier, 1917 (non consulté).
\item \textsc{Kipling}, Rudyard, \textit{The Irish Guards in the Great War}, vol.~1, Londres, Macmillan, 1923. Project Gutenberg n\textsuperscript{o}~64638.
\item \textsc{Laboureyras}, P., \textit{La Destruction d'une cité picarde et d'une basilique mariale. La ville d'Albert avant et pendant la guerre, 1914--1915}, 1916, rééd. 2012 (cité de seconde main).
\item \textsc{Lafon}, Gabriel, \textit{Les Chemins de fer français pendant la guerre}, thèse pour le doctorat, Université de Paris, Faculté de droit, Paris, A.~Rousseau, 1922. \url{https://archive.org/details/IA41551422_0037}.
\item \textsc{Le Hénaff} (colonel) et \textsc{Bornecque}, Henri (capitaine), \textit{Les Chemins de fer français et la guerre}, Paris, Chapelot, 1922 (Gallica, \url{https://gallica.bnf.fr/ark:/12148/bpt6k65294286} ; non consulté).
\item \textsc{Linfoot}, Arthur, journal de guerre (58th Field Ambulance, RAMC), édition annotée en ligne, \url{https://arthurlinfoot.org.uk}.
\item \textsc{Maeterlinck}, Maurice, \textit{Les Débris de la guerre}, Paris, Eugène Fasquelle, Bibliothèque-Charpentier, 1916 (exemplaire consulté : 16\textsuperscript{e}~mille, 1918). \url{https://archive.org/details/lesdbrisdelagu00maetuoft}.
\item \textsc{March}, Lucien, \textit{Mouvement des prix et des salaires pendant la guerre}, Paris, Presses universitaires de France, et New Haven, Yale University Press, 1925 (Publications de la Dotation Carnegie pour la paix internationale). \url{https://archive.org/details/fp-0269}.
\item \textsc{Marie-Antoine} (R.~P.), \textit{Le Livre d'or des proscrits}, 12\textsuperscript{e}~éd., Paris, V.~Palmé, et Toulouse, É.~Privat, 1881. \url{https://archive.org/details/lelivredordespr00antogoog}.
\item \textsc{Martin}, Alexis, \textit{Une visite à Amiens. Historique. La ville haute, la ville basse, les faubourgs. Hommes célèbres nés en Picardie}, Paris, A.~Hennuyer, 1896 (BnF, notice de la réédition, La Vague verte, 2016, ark:/12148/cb450737666).
\item \textsc{Maurois}, André, \textit{Les Silences du colonel Bramble}, Paris, Grasset, 1918 (\url{https://archive.org/details/lessilencesducol00maur}) ; \textit{The Silence of Colonel Bramble}, trad. Thurfrida Wake, Londres et New York, John Lane, 1920 («~Biographical Introduction~»).
\item \textsc{Maze}, Paul, \textit{A Frenchman in Khaki}, Londres, Heinemann, 1934 (cité d'après Cowley 2016 ; non consulté).
\item \textsc{Michelin \& Cie}, \textit{Amiens before and during the War}, Clermont-Ferrand, Michelin, 1919. Project Gutenberg n\textsuperscript{o}~48347.
\item \textsc{Michelin \& Cie}, \textit{Illustrated Michelin Guides to the Battlefields (1914--1918). The Somme}, vol.~1 : \textit{The First Battle of the Somme (1916--1917)}, Clermont-Ferrand, Michelin, 1919. Project Gutenberg n\textsuperscript{o}~49122 ; \url{https://archive.org/details/cu31924088058940}.
\item \textsc{Michelin \& Cie}, \textit{St-Quentin -- Cambrai, la ligne Hindenburg} (Guides illustrés Michelin des champs de bataille), Clermont-Ferrand, Michelin, 1921. \url{https://archive.org/details/stquentincambrai00cler}.
\item \textsc{Murphy}, C.~C.~R., \textit{The History of the Suffolk Regiment, 1914--1927}, Londres, Hutchinson, [1928]. \url{https://archive.org/details/historysuffolkregimentmurphy}.
\item \textsc{Nichols}, G.~H.~F. («~Quex~»), \textit{The 18th Division in the Great War}, Édimbourg et Londres, W.~Blackwood \& Sons, 1922. \url{https://archive.org/details/18thdivision-greatwar}.
\item \textsc{O'Neill}, H.~C., \textit{The Royal Fusiliers in the Great War}, Londres, Heinemann, 1922. \url{https://archive.org/details/royalfusiliersin00onei}.
\item \textsc{Pearse}, H.~W., et \textsc{Sloman}, H.~S., \textit{History of the East Surrey Regiment}, vol.~II (1914--1917), Londres, The Medici Society, 1923. \url{https://archive.org/details/histeastsurreyregimentvol2}.
\item \textsc{Petre}, F.~Loraine, \textit{The History of the Norfolk Regiment 1685--1918}, vol.~II, Norwich, Jarrold \& Sons, [1925]. \url{https://archive.org/details/norfolkregimentvol2}.
\item \textsc{Sassoon}, Siegfried, \textit{Memoirs of an Infantry Officer}, Londres, Faber \& Faber, 1930.
\item \textsc{Shakespear}, J., \textit{The Thirty-Fourth Division, 1915--1919. The Story of its Career from Ripon to the Rhine}, Londres, H.~F.~\& G.~Witherby, 1921. \url{https://archive.org/details/thirtyfourthdivision19151919}.
\item \textsc{Shaw Sparrow}, W., \textit{The Fifth Army in March 1918}, introduction du général Sir H.~Gough, Londres, John Lane, 1921 (consulté sans emprunt).
\item \textsc{Whittem}, Arthur F., et \textsc{Long}, Percy W., \textit{French for Soldiers}, Cambridge (Mass.), Harvard University Press, 1917.
\item \textsc{Wyrall}, Everard, \textit{The Die-Hards in the Great War. A History of the Duke of Cambridge's Own (Middlesex Regiment), 1914--1919}, vol.~I, Londres, Harrison \& Sons, [1926]. \url{https://archive.org/details/diehardsvol1}.
\item \textsc{Yapp}, Sir Arthur K., \textit{The Romance of the Red Triangle. The Story of the Coming of the Red Triangle and the Service Rendered by the Y.M.C.A. to the Sailors and Soldiers of the British Empire}, Londres, Hodder \& Stoughton, 1918. Project Gutenberg n\textsuperscript{o}~32998.
\item \textsc{Zyromski}, Ernest, \textit{L'Orgueil humain}, Paris, A.~Colin, 1904.
\item \textsc{Zyromski}, Ernest, \textit{Maurice de Guérin}, Paris, Librairie Armand Colin, 1921. \url{https://archive.org/details/mauricedegurin00zyrouoft}.
\end{biblist}

\subsubsection*{C. Dictionnaires et lexiques}

\begin{biblist}
\item \textsc{Block}, Maurice, \textit{Dictionnaire de l'administration française}, 5\textsuperscript{e}~éd. refondue sous la direction d'Édouard Maguéro, Paris et Nancy, Berger-Levrault, 1905. \url{https://archive.org/details/dictionnairedela01bloc}.
\item \textsc{Clifton}, C.~E., et \textsc{Grimaux}, A., \textit{Nouveau dictionnaire anglais-français et français-anglais}, t.~I, Paris, Garnier, 1883 (exemplaire consulté : Montréal, Beauchemin). \url{https://archive.org/details/nouveaudictionna01clif}.
\item \textit{Dictionnaire des termes militaires et de l'argot poilu}, Paris, Librairie Larousse, [1916]. \url{https://archive.org/details/dictionnairedest00pariuoft}.
\item \textsc{Fraser}, Edward, et \textsc{Gibbons}, John, \textit{Soldier and Sailor Words and Phrases}, Londres, Routledge, 1925. \url{https://archive.org/details/bwb_Y0-BCE-327}.
\item \textit{French-English Dictionary of Military Terms and of Other Terms Useful to Officers}, Londres, HMSO, 1919. \url{https://archive.org/details/bwb_W9-CFM-565}.
\item \textit{The Oxford English Dictionary}, Oxford, 1933.
\item \textsc{Partridge}, Eric, \textit{A Dictionary of Slang and Unconventional English}, Londres, Routledge, 1937.
\end{biblist}

\subsubsection*{D. Presse et périodiques}

\begin{biblist}
\item \textsc{Anonyme} («~A Divisional Claims Officer~»), «~Claims~», \textit{Blackwood's Magazine}, vol.~201, n\textsuperscript{o}~1220, juin 1917, p.~933--941. \url{https://archive.org/details/blackwoodsmagazi201edinuoft}.
\item \textsc{Barbé}, G., «~L'hiver 1916-1917~», \textit{Journal d'agriculture pratique}, nouvelle série, t.~30 (81\textsuperscript{e}~année), Paris, Librairie agricole de la Maison rustique, 1917, p.~181--182. \url{https://archive.org/details/journaldagricult81pari}.
\item \textit{Bulletin de la Société de géographie de Toulouse}, 1\textsuperscript{er}~avril 1941 et 1\textsuperscript{er}~juin 1942 (Rosalis, ark:/12148/bpt6k6282699q et bpt6k6369236p).
\item \textit{Bulletin des armées de la République}, 2\textsuperscript{e}~année, n\textsuperscript{o}~186, 22~mars 1916, p.~2, «~Les permissions~» ; 4\textsuperscript{e}~année, n\textsuperscript{o}~268, 17~octobre 1917, p.~2 (La contemporaine, Argonnaute).
\item \textit{La Dépêche} (Toulouse) : 27~juin 1914, p.~4 ; 29~août 1916, p.~1 ; 28~janvier 1917, p.~2 ; 30~janvier 1917, p.~3 ; 31~janvier 1917, p.~2--3 ; 2~février 1917, p.~3--4 («~Neiges et glaces. Note de la mairie~») ; 3 et 4~février 1917, p.~3 ; 8~août 1917, p.~3 ; 15~mai 1919, p.~3 ; 10~juillet 1922, p.~4 (Rosalis et Gallica ; ark:/12148/bpt6k41508539, bpt6k4150549c, bpt6k4146920v, bpt6k4146922p, bpt6k41469233, bpt6k4146925x, bpt6k4146926b, bpt6k4146927r, bpt6k4147111f, bpt6k4147757h, bpt6k41418785).
\item \textit{L'Auta}, 1\textsuperscript{er}~juin 1919, p.~12 (Rosalis, ark:/12148/bpt6k96057467) ; janvier 2002 : voir \textsc{Daraux} 2002 (III).
\item \textit{L'Express du Midi} (Toulouse) : 11~février 1907, p.~2 ; 25~décembre 1911, p.~3 ; 18 et 19~septembre 1914, p.~2 ; 10~janvier 1916, p.~3 ; 17, 28 et 31~janvier 1917, p.~3 ; 1\textsuperscript{er}~février 1917, p.~3 («~Cours et conférences~» ; «~Maurice de Guérin à la Chênaie~») ; 3~février 1917, p.~4 ; 9~septembre 1923, p.~3 ; 23~décembre 1926, p.~5 (faire-part d'Émile Crouzet) ; 27 et 28~août 1932, p.~4--5 (Rosalis ; ark:/12148/bpt6k53687164, bpt6k5383179k, bpt6k53841631, bpt6k5384164f, bpt6k5384640r, bpt6k5385012p, bpt6k5385023g, bpt6k5385026q, bpt6k53850274, bpt6k5385029z, bpt6k53874222, bpt6k5388614h, bpt6k5391104h, bpt6k5391105x).
\item \textit{Le Midi socialiste} (Toulouse) : 31~mars 1911, p.~3 ; 5~février 1917, p.~2 ; 14~avril 1917, p.~2 ; 5~décembre 1922, p.~3 ; 17~juillet 1923, p.~3 ; 13~mai 1941, p.~2 (Rosalis ; ark:/12148/bpt6k53435029, bpt6k5346061c, bpt6k53461283, bpt6k5348099f, bpt6k5348317m, bpt6k5354757j).
\item \textit{Le Télégramme} (Toulouse) : 31~mars 1911, p.~4 ; 23~novembre 1912, p.~3 ; 27~janvier 1913, p.~4 ; 27~juin 1914, p.~4 ; 5~février 1915, p.~2 ; 10~mars 1916, p.~4 ; 27~janvier 1917, p.~2 ; 30~janvier 1917, p.~3 ; 1\textsuperscript{er} et 4~février 1917, p.~3 ; 15~avril 1917, p.~3 ; 3~juillet 1917, p.~2 ; 1\textsuperscript{er}~août 1925, p.~6 ; 24~décembre 1926, p.~4 (Rosalis ; ark:/12148/bd6t53864395x, bd6t538649898, bd6t538650548, bd6t538655703, bd6t53865791z, bd6t53866190s, bd6t53866512z, bd6t538665156, bd6t538665171, bd6t53866520h, bd6t538665905, bd6t53866670r, bd6t54833338z, bd6t548408062).
\item \textit{Les Langues modernes}, bulletin de l'Association des professeurs de langues vivantes de l'enseignement public, XVI, 1918, p.~5--6 et 70 ; XVII, 1919, n\textsuperscript{o}~2, p.~102--105, et n\textsuperscript{o}~3, «~Chronique du mois~», p.~184 ; XVIII, 1920, p.~86 et 305 ; XIX, 1921, p.~339. \url{https://archive.org/details/languesmodernes16asso} ; \url{https://archive.org/details/languesmodernes18asso}.
\item \textit{Recueil de l'Académie de Montauban}, année 1932, p.~47--57 (H.~\textsc{Combes}, «~Un «~Poilu~» de la Révolution : le sergent Landou de Montauban~», lu le 4~janvier 1932) ; 2\textsuperscript{e}~série, t.~XLVIII, année 1938, Montauban, Forestié, 1939 (H.~\textsc{Pouvillon}, discours de réception du 7~février 1938, p.~150--151, et discours du président Caldesaigues, p.~147) (Gallica, ark:/12148/bpt6k65323798 ; copies au dépôt).
\item \textit{Revue de l'enseignement des langues vivantes}, 35, 1918. \url{https://archive.org/details/revuedelenseigne35pariuoft}.
\item \textit{The London Gazette}, suppléments : n\textsuperscript{o}~29235, 19~juillet 1915, p.~7057 ; n\textsuperscript{o}~29321, 9~octobre 1915, p.~9989--9990 ; n\textsuperscript{o}~29403, 15~décembre 1915, p.~12533 ; n\textsuperscript{o}~29585, 17~mai 1916, p.~4941 ; n\textsuperscript{o}~29787, 18~octobre 1916, p.~10039 ; n\textsuperscript{o}~30111, 4~juin 1917, p.~5468 et 5470 ; n\textsuperscript{o}~30716, 3~juin 1918, p.~6457--6459 ; n\textsuperscript{o}~31741, 19~janvier 1920, p.~823 ; n\textsuperscript{o}~32343, 2~juin 1921, p.~4387. \url{https://www.thegazette.co.uk}.
\end{biblist}

%======================================================================
\subsection*{III. Études et instruments de recherche}

\begin{biblist}
\item \textsc{Australian War Memorial}, \textit{Backgrounder : The Battle of Fromelles, 19--20 July 1916}, PDF, AWM.
\item \textsc{Baynes}, John, \textit{Far from a Donkey : The Life of General Sir Ivor Maxse}, Londres, Brassey's, 1995 (non consulté).
\item \textsc{Bazin}, Jacques, \textit{Toulouse : l'histoire du réseau}, 2008 (cité d'après Wikipédia).
\item \textsc{Beach}, Jim, et \textsc{Bruce}, James, «~Improvising Language Capability : The British Army's Corps of Interpreters, 1914--1915~», \textit{War in History}, 2023, DOI 10.1177/09683445221140305 ; manuscrit accepté : \url{https://nectar.northampton.ac.uk/id/eprint/18436/}.
\item \textsc{Beadon}, R.~H., \textit{The Royal Army Service Corps : A History of Transport and Supply in the British Army}, vol.~II, Cambridge, Cambridge University Press, 1931 (non consulté).
\item \textsc{Cowley}, Peter, «~Declining the interpreter's role in World War~I~», \textit{Linguistica Antverpiensia, New Series -- Themes in Translation Studies}, 15, 2016, p.~72--88. \url{https://lans-tts.uantwerpen.be/index.php/LANS-TTS/article/view/394}.
\item \textsc{Cronier}, Emmanuelle, \textit{Permissions et permissionnaires du front à Paris pendant la Grande Guerre}, thèse, Université Paris~1, 2005.
\item \textsc{Cronier}, Emmanuelle, \textit{Permissionnaires dans la Grande Guerre}, Paris, Belin, 2013 (cité de seconde main, d'après Wikipédia).
\item \textsc{Daraux}, Jean-Pierre, «~L'œuvre de Louis Mazetier (1888--1952) au monastère de Blagnac~», \textit{L'Auta}, janvier 2002 (Rosalis, ark:/12148/bpt6k9610682g).
\item \textsc{Darmon}, Pierre, \textit{Vivre à Paris pendant la Grande Guerre}, Paris, Fayard, 2002 (non consulté).
\item \textsc{Anonyme}, «~L'église du Jésus à Toulouse : architecture et décor~», \textit{Mémoires de la Société archéologique du Midi de la France}, t.~LX, 2000, p.~193--210. \url{https://societearcheologiquedumidi.fr/_samf/memoires/t_60/193-210N.PDF}.
\item \textsc{Espel}, Arnaud, \textit{Fonds de l'Office notarial de Toutencourt, aujourd'hui Acheux-en-Amiénois (1729--1904). Répertoire méthodique détaillé}, sous la dir. d'Anne Lejeune, Amiens, Archives départementales de la Somme, 2018. \url{https://archives.somme.fr/document/toutencourt}.
\item \textsc{Farebrother}, Martin J.~B. et Joan, \textit{Allied Railways of the Western Front. Narrow Gauge in the Somme Sector. Before, During and After the First World War}, Barnsley, Pen \& Sword Transport, 2018 (cité d'après Wikipédia, qui le date de 2016).
\item \textsc{Feller-Vest}, Veronika, «~Commynes, Philippe de~», \textit{Dictionnaire historique de la Suisse}, version en ligne du 18~novembre 2015. \url{https://hls-dhs-dss.ch/fr/articles/018723/}.
\item \textsc{Flutre}, Louis-Fernand, \textit{Le Parler picard de Mesnil-Martinsart (Somme). Phonétique, morphologie, syntaxe, vocabulaire}, Genève, Droz, et Lille, Giard, 1955 (BnF, \url{https://catalogue.bnf.fr/ark:/12148/cb321136530}) ; \textit{Mesnil-Martinsart (Somme), essai d'histoire locale}, Lyon, 1955.
\item \textsc{Gibson}, Craig, «~The British Army, French Farmers and the War on the Western Front 1914--1918~», \textit{Past \& Present}, 180, 2003 ; \textit{Behind the Front : British Soldiers and French Civilians, 1914--1918}, Cambridge, Cambridge University Press, 2014 (non consultés ; voir \textsc{Lukasik} 2016).
\item \textsc{Greenhalgh}, Elizabeth, \textit{Victory through Coalition : Britain and France during the First World War}, Cambridge, Cambridge University Press, 2005 (non consulté).
\item \textsc{Guillot}, Hélène, et \textsc{Rohfrischt}, Anne, \textit{Missions militaires françaises (archives rapatriées de Russie). Sous-série GR~17~NN, 1916--1922. Répertoire numérique détaillé}, Vincennes, Service historique de la Défense, 2015. \url{https://www.servicehistorique.sga.defense.gouv.fr/sites/default/files/notices_files/SHDGR_REP_17NN_MMF.pdf}.
\item \textsc{Heimburger}, Franziska, «~Fighting Together : Language Issues in the Military Coordination of First World War Allied Coalition Warfare~», dans H.~Footitt et M.~Kelly (dir.), \textit{Languages and the Military : Alliances, Occupation and Peace Building}, Londres, Palgrave Macmillan, 2012, p.~49--55 (non consulté).
\item \textsc{Heimburger}, Franziska, «~Imagining Coalition Warfare ? French and British Military Language Policy before 1914~», \textit{Francia}, 40, 2013, p.~399--403 (non consulté).
\item \textsc{Heimburger}, Franziska, «~Of Go-Betweens and Gatekeepers : Considering Disciplinary Biases in Interpreting History through Exemplary Metaphors. Military Interpreters in the Allied Coalition during the First World War~», dans B.~Fischer et M.~Nisbeth Jensen (dir.), \textit{Translation and the Reconfiguration of Power Relations}, Berlin, Lit Verlag, 2013 (non consulté).
\item \textsc{Heimburger}, Franziska, \textit{Mésentente cordiale ? Langues et coalition alliée sur le front de l'ouest de la Grande Guerre}, thèse de doctorat, EHESS, 2014 (non consultée).
\item \textit{International Directory of Company Histories}, vol.~19, St.~James Press, 1998, «~Daily Mail and General Trust plc History~» (\url{https://www.fundinguniverse.com}).
\item \textsc{Lacombe}, Christian, et \textsc{Malterre}, Patrice, \textit{Omnibus, tramways et autobus de Toulouse}, Toulouse, éd. du Cabri, 1983 (cité d'après Wikipédia).
\item \textsc{Lukasik}, S.~H., compte rendu de C.~Gibson, \textit{Behind the Front}, \textit{Michigan War Studies Review}, 2016-052. \url{https://michiganwarstudiesreview.com/2016-052.htm}.
\item \textsc{Neiberg}, M.~S., compte rendu de E.~Greenhalgh (éd.), \textit{Liaison}, \textit{British Journal for Military History}, 4/1, novembre 2017, p.~124--125.
\item \textsc{Nicot}, Jean, \textsc{Schillinger}, Philippe, et \textsc{Ficat}, Josette, \textit{Inventaire sommaire des archives de la guerre 1914--1918. Grand quartier général des armées alliées (15~N), GQG (16~N), Missions militaires françaises (17~N)}, Vincennes, Service historique de l'Armée, 1972.
\item \textsc{Offenstadt}, Nicolas (dir.), \textit{Le Chemin des Dames. De l'événement à la mémoire}, Paris, Stock, 2004.
\item \textsc{Pidgeon}, Trevor, \textit{Boom Ravine}, coll. «~Battleground Europe~», Barnsley, Leo Cooper / Pen \& Sword, 1998 (non consulté).
\item \textsc{Séchet}, Guillaume, «~Le supplice du terrible mois de février 1917~», Météo-Paris. \url{https://www.meteo-paris.com/actualites/le-supplice-du-terrible-mois-de-fevrier-1917}.
\item \textsc{Shiman}, Lilian Lewis, «~The Band of Hope Movement : Respectable Recreation for Working-Class Children~», \textit{Victorian Studies}, 17/1, 1973, p.~49--74.
\item \textsc{Sutra}, Romy, \textit{«~La loi à la main~». Militantisme juridique et défense religieuse au temps de l'affirmation de la République. L'action du comité de jurisconsultes des congrégations (1880--1905)}, Toulouse, Presses de l'Université Toulouse Capitole, chap.~2, «~La résistance judiciaire aux décrets du 29~mars 1880~». \url{https://books.openedition.org/putc/14400}.
\item \textsc{Tattersfield}, David, «~The workhorse of the staff : the Vauxhall D-Type staff car~», Western Front Association, 9~janvier 2021. \url{https://www.westernfrontassociation.com/world-war-i-articles/the-workhorse-of-the-staff-the-vauxhall-d-type-staff-car}.
\item \textsc{Thompson}, John Handby, \textit{The Free Church Army Chaplain 1830--1930}, thèse de doctorat, Université de Sheffield, 1990 (résumé). \url{https://etheses.whiterose.ac.uk/id/eprint/1785/}.
\item \textsc{Ward}, Graham, «~Wren's Shoe Polish and Northampton~», \textit{Miscellanea Edintone}, 20~juin 2022, mis à jour le 2~décembre 2024. \url{https://edintone.com/wrens-shoe-polish-and-northampton}.
\item \textsc{Zeller}, Gaston, «~Un faux du XVII\textsuperscript{e} siècle. L'édit de Louis~XI sur la poste~», \textit{Revue historique}, t.~180, 1937 (cité d'après Zeller 1946).
\item \textsc{Zeller}, Gaston, «~Procès à reviser ? Louis~XI, la noblesse et la marchandise~», \textit{Annales. Économies, Sociétés, Civilisations}, 1, n\textsuperscript{o}~4, 1946, p.~331--341. \url{https://www.persee.fr/doc/ahess_0395-2649_1946_num_1_4_3239}.
\end{biblist}

%======================================================================
\subsection*{IV. Bases de données et sites}

\begin{biblist}
\item \textbf{Archive.org} (Internet Archive) et \textbf{Project Gutenberg} : exemplaires numérisés cités aux parties II et III.
\item \textbf{Archives municipales de Toulouse}, «~Histoire des marchés toulousains~» et «~Marchés dans la ville~», \url{https://archives.mairie-toulouse.fr}.
\item \textbf{Australian War Memorial}, REL C1245196, «~Fundraising badges «~Church Army Huts. Hut Day~»~». \url{https://www.awm.gov.au/collection/C1245196}.
\item \textbf{Base adresse nationale} (\url{https://api-adresse.data.gouv.fr/search/}) et \textbf{API Découpage administratif} (\url{https://geo.api.gouv.fr/communes}) : coordonnées des mairies ; distances à vol d'oiseau.
\item \textbf{Bedfordshire Archives}, catalogue en ligne : voir I\,B.
\item \textbf{BnF}, catalogue général : «~[Miroir des modes]~», Paris, Butterick, 1902--1934, ark:/12148/cb34387442t ; «~Le Miroir des modes~», 1897--1934, ark:/12148/cb32816734x ; \textbf{data.bnf.fr}, notice d'autorité «~Ernest Zyromski (1862--1933)~», \url{https://data.bnf.fr/12976861}.
\item \textbf{Calendar-12}, «~1916 Moon Phases~» et «~Moon Phases 1917~» (\url{https://calendar-12.com/moon_phases/1916}, \url{https://calendar-12.com/moon_phases/1917}) ; \textbf{iCalendrier}, «~Calendrier lunaire février 1917~» (\url{https://icalendrier.fr/calendrier-lunaire/fevrier-1917}) ; \textbf{TheSkyLive}, «~Moon Calendar 1917~» (\url{https://theskylive.com/moon-calendar?year=1917}).
\item \textbf{CanalPlan AC}, «~Lys~» (écluse~5 et barrage de Bac-Saint-Maur), \url{https://canalplan.co.uk/waterway/zuwi}.
\item \textbf{Château des ducs de Bretagne -- Musée d'histoire de Nantes}, «~LU, a Nantes story~», exposition en ligne, Google Arts \& Culture. \url{https://artsandculture.google.com/exhibit/lu-a-nantes-story/YwJSmyHH3im9Kg}.
\item \textbf{Church Army}, «~Church Army Through the World Wars~». \url{https://churcharmy.org/be-inspired/catalystjuly22/ca-through-the-world-wars/}.
\item \textbf{Commonwealth War Graves Commission} [CWGC], notices historiques des cimetières (\url{https://www.cwgc.org/visit-us/find-cemeteries-memorials/cemetery-details/}, suivi du numéro) : Acheux British (11501) ; Agny (24503) ; Albert Communal Cemetery Extension (61300) ; Anzac Cemetery, Sailly (18802) ; Artillery Wood, Boezinge (9700) ; Aveluy Communal Cemetery Extension (18700) ; Bailleul Communal Cemetery Extension (2000048) ; Beauval Communal (3601) ; Bertrancourt Military (11502) ; Bienvillers (26502) ; Bois-Grenier (63405) ; Bouzincourt Communal Cemetery Extension (2000051) ; Brown's Road, Festubert (2000035) ; Cabaret-Rouge, Souchez (64600) ; Cambrin Churchyard Extension (16800) ; Cambrin Military (3400) ; Chapelle-d'Armentières Old Military (63402) ; Cité Bonjean, Armentières (27801) ; Connaught, Thiepval (58600) ; Contay British (8000) ; Courcelette British (2000045) ; Croix-du-Bac British, Steenwerck (28202) ; Doullens Communal Cemetery Extension No.~1 (8101) ; Dud Corner, Loos (58700) ; Englebelmer Communal Cemetery Extension (4003798) ; Erquinghem-Lys Churchyard Extension (2000041) ; Estaires Communal Cemetery and Extension (26801) ; Faubourg d'Amiens, Arras (28800) ; Foncquevillers Military (20200) ; Forceville Communal Cemetery and Extension (11400) ; Frankfurt Trench British, Beaumont-Hamel (2000153) ; Fromelles (Pheasant Wood) (5001073) ; Gézaincourt Communal Cemetery Extension (3604) ; Gommecourt British No.~2, Hébuterne (65200) ; Gommecourt Wood New, Foncquevillers (31400) ; Gordon, Mametz (59803) ; Gordon Dump, Ovillers-La Boisselle ; Grandcourt Road (60601) ; Guards Cemetery, Windy Corner, Cuinchy (4003815) ; Hazebrouck Communal (4700) ; Hébuterne Military (33301) ; Hédauville Communal Cemetery Extension (28601) ; Hooge Crater (52700) ; Houplines (28401) ; Kemmel Chateau (15401) ; Lancashire Cottage (51600) ; Laventie (34300) ; Le Grand Beaumart British, Steenwerck (20802) ; Le Touquet Railway Crossing (8803) ; Lille Southern (28001) ; Lonsdale, Authuille (59400) ; Louvencourt Military (11500) ; Mailly-Maillet Communal Cemetery Extension (28602) ; Maroc, Grenay (58200) ; Martinsart British (18100) ; Méaulte Military ; Mesnil Communal Cemetery Extension (18101) ; Messines Ridge British (55400) ; Mill Road, Thiepval (60600) ; Mont Noir Military (27403) ; Neuve-Chapelle (66005) ; New Munich Trench British, Beaumont-Hamel (30900) ; Ploegsteert Wood (51602) ; Puchevillers British (57401) ; Railway Dugouts (56000) ; Regina Trench, Grandcourt (23000) ; Rue-Pétillon (30400) ; Sailly-sur-la-Lys Churchyard (18800) ; Sanctuary Wood (54700) ; Serre Road No.~2 (67200) ; St.~Vaast Post, Richebourg-l'Avoué (63000) ; Steenwerck Communal (28201) ; Stump Road, Grandcourt (30901) ; Trois Arbres, Steenwerck (2000052) ; Varennes Military (7200) ; V.C. Corner (78900) ; Warlincourt Halte British, Saulty ; Warloy-Baillon Communal Cemetery Extension (7301) ; Wytschaete Military (51700) ; Zouave Valley, Souchez (62202). Fiche individuelle : Second Lieutenant David Arnold Roeber, 3rd Bn attd 7th Bn Bedfordshire Regiment, mort le 14~août 1916, Ration Farm Military Cemetery, La Chapelle-d'Armentières, II.C.2. \textbf{CWGC Foundation}, «~A Guide's Guide to the Somme : Serre Road Cemetery No.~2~», 23~août 2022.
\item \textbf{Communauté de communes Millau Grands Causses}, «~Un territoire attractif~» (brochure PDF).
\item \textbf{DBnomics}, série INSEE COEFF-EURO-FRANC (base 2015), et \textbf{INSEE}, convertisseur franc-euro (\url{https://www.insee.fr/fr/information/2417794}).
\item \textbf{Des villages de Cassini aux communes d'aujourd'hui} (EHESS), Saint-Pol-sur-Ternoise, recensement de 1911 (consulté par l'intermédiaire de Wikipédia).
\item \textbf{Diocèse d'Amiens}, «~Notre-Dame de Brebières ou la «~Vierge penchée~»~», 13~mars 2018. \url{https://www.amiens.catholique.fr}.
\item \textbf{Églises ouvertes}, notice «~Église Saint-Vaast, Bailleul~» ; \textbf{Cirkwi}, points d'intérêt 4393487 (Bailleul) et 2443051 (Ville d'Armentières, circuit du patrimoine, «~Richesses architecturales : l'église Saint-Vaast, la reconstruction~»).
\item \textbf{Espenak}, Fred, \textit{Six Millennium Catalog of Phases of the Moon}, table «~Phases of the Moon : 1901 to 2000~», \url{https://astropixels.com/ephemeris/phasescat/phases1901.html} ; positions astronomiques calculées par l'éditeur avec PyEphem~4.2.1.
\item \textbf{Gallica} (BnF) et \textbf{Rosalis} (Bibliothèque de Toulouse, \url{https://rosalis.bibliotheque.toulouse.fr/ark:/12148/}, suivi de l'identifiant) : périodiques cités en II\,D.
\item \textbf{Geneawiki}, «~12145 -- Millau~». \url{https://fr.geneawiki.com/wiki/12145_-_Millau}.
\item \textbf{Grace's Guide to British Industrial History} : «~F.~C.~Calvert and Co~», «~Nugget Polish Co~», «~Sunbeam Motor Car Co~», «~Wolseley Motors~». \url{https://www.gracesguide.co.uk}.
\item \textbf{Hérodote.net}, «~19~juillet 1900 : inauguration du métro de Paris~». \url{https://www.herodote.net}.
\item \textbf{IGN}, Géoplateforme, service de géocodage, index des points d'intérêt (BD TOPO) : lieu-dit habité «~le Quesnoy~», Puchevillers. \url{https://data.geopf.fr/geocodage/search}.
\item \textbf{Imperial War Museum}, photographie Q~961 (Lt Ernest Brooks), 9~août 1916 : le roi à la Third Army Trench Mortar School de Ligny-Saint-Flochel (Ministry of Information First World War Official Collection). \url{https://www.iwm.org.uk/collections/item/object/205072383}.
\item \textbf{INSEE}, fichier des personnes décédées, interrogé par l'intermédiaire de matchID (\url{https://deces.matchid.io}) ; répertoire SIRENE (\url{https://annuaire-entreprises.data.gouv.fr}), établissements Barbry, Sailly-sur-la-Lys.
\item \textbf{Inventaire général du patrimoine culturel}, Région Hauts-de-France (\url{https://inventaire.hautsdefrance.fr}) : dossiers IA80009826 («~Hameau de Barlette~», Ribeaucourt, 2010), IA80000292 (château de Beaucourt-sur-l'Hallue, 2000), «~Le village de Franqueville~», «~L'ancienne ligne de chemin de fer de Frévent à Gamaches~», usine de blanchiment Delecaille dite Blanchisserie linière de France (Erquinghem-Lys) ; Région Auvergne-Rhône-Alpes (\url{https://patrimoine.auvergnerhonealpes.fr}) : «~Ligne Brioude -- Saint-Flour~», «~Ligne Neussargues -- Loubaresse -- Béziers~», «~Gare de Neussargues~».
\item «~La gravure sur nacre, cet art oublié des bagnards~», \textit{Les Nouvelles calédoniennes}. \url{https://www.lnc.nc/article/la-gravure-sur-nacre-cet-art-oublie-des-bagnards}.
\item \textbf{London Gazette} (\url{https://www.thegazette.co.uk}) et \textbf{TNA Discovery} (\url{https://discovery.nationalarchives.gov.uk}) : voir II\,D et I\,B.
\item \textbf{Masonic Great War Project}. \url{https://www.masonicgreatwarproject.org.uk}.
\item \textbf{Mémoire des hommes} (ministère des Armées) et \textbf{Léonore} : non interrogés (accès impossible depuis l'environnement de travail).
\item \textbf{Ministère de la Culture, base Mérimée} (Plateforme ouverte du patrimoine, \url{https://www.pop.culture.gouv.fr} ; jeu de données «~Liste des immeubles protégés au titre des monuments historiques~», \url{https://data.culture.gouv.fr} ; relais \url{https://monumentum.fr}) : IA00062810 (blanchisserie Auguste Mahieu, Erquinghem-Lys) ; IA00076408 (sucrerie dite Cie sucrière du Canton d'Acheux, puis Normand et Cie, Acheux-en-Amiénois) ; PA00086489 (église Saint-Vincent-de-Paul, Paris~X\textsuperscript{e}) ; PA00094498 (cathédrale Saint-Étienne, Toulouse) ; PA00094499 (chapelle Notre-Dame-de-Nazareth, Toulouse) ; PA00094527 (gare de Toulouse-Matabiau) ; PA00107358 (beffroi et hôtel de ville, Bailleul) ; PA00107549 (église Saint-Éloi, Hazebrouck) ; PA00107783 (château dit de Philippe de Comines, Renescure) ; PA00107784 (château de Zuthove, Renescure) ; PA00107816 (église Saint-Érasme, Sercus) ; PA00107832 ; PA00108287 (église Saint-Hilaire, Frévent) ; PA00108332 (église Saint-Flochel, clocher, Ligny-Saint-Flochel) ; PA00108388 ; PA00108423 ; PA00108424 (église, Saint-Pol-sur-Ternoise) ; PA00108425 ; PA00116138 (ancienne église Saint-Pierre, Doullens) ; PA00116177 ; PA00116193 (beffroi, Lucheux) ; PA00116194 (château, Lucheux) ; PA00116195 (église Saint-Léger, Lucheux) ; PA00116224 (château, Raincheval) ; PA31000123 (musée Georges-Labit, Toulouse) ; PA59000059 (ancienne usine de blanchiment Mahieu, Erquinghem-Lys) ; PA59000085 ; PA59000231 (église Saint-Jean-Baptiste, Steenwerck) ; PA59000232 (monument aux morts, Bailleul) ; PA80000002 (basilique Notre-Dame de Brebières, Albert) ; PA80000004 (trou de mine de La Boisselle) ; PA80000063 (jardin public, Albert) ; PA80000086 (usine Saint Frères, Beauval). Aucune notice pour Estaires ni pour Saint-Michel-sur-Ternoise (interrogation du 18~octobre 2026).
\item \textbf{Ministère de la Culture, base Palissy} (POP ; export \texttt{palissy.csv} du 16~octobre 2026) : PM59000995 (jubé, église Saint-Vaast, Lynde) ; PM59003991 et PM59003992 (Steenwerck, fonts baptismaux ; chapiteau et fût de colonne) ; PM59007232 (F.~Ch.~Baude, \textit{Vue de l'église Saint-Vaast d'Armentières après les combats de la Première Guerre mondiale}, 1919, hôtel de ville d'Armentières) ; PM62001934 (orgue, Saint-Pol-sur-Ternoise) ; PM80005288 (Vierge à l'Enfant, église Saint-Martin, Hérissart) ; objets mobiliers des communes d'Erquinghem-Lys, Fleurbaix, Hazebrouck, Renescure (église Notre-Dame-de-l'Assomption), Sailly-sur-la-Lys, Saint-Jans-Cappel et Sercus.
\item \textbf{Musée d'Orsay}, «~Histoire du musée~». \url{https://www.musee-orsay.fr/fr/articles/histoire-du-musee-60}.
\item \textbf{Office de tourisme de Toulouse} (\url{https://www.toulouse-tourisme.com}) : «~Amphithéâtre romain de Purpan et thermes d'Ancely~» ; «~Marché couvert Victor Hugo~» ; «~Notre-Dame de la Dalbade~».
\item \textbf{Office québécois de la langue française}, \textit{Grand Dictionnaire terminologique}, fiche «~pot au noir~».
\item \textbf{Online Etymology Dictionary}, «~Archibald~» (\url{https://www.etymonline.com/word/Archibald}) ; \textbf{The Online Books Page}, «~Martin, Alexis, 1834--1904~» (notice HathiTrust 009016325).
\item \textbf{OpenAlex}, notice de la thèse de F.~Heimburger (DOI 10.4000/acrh.6310).
\item \textbf{OpenStreetMap}, service Nominatim (\url{https://nominatim.openstreetmap.org/}) : géocodage des localités ; distances à vol d'oiseau entre chefs-lieux.
\item \textbf{Ordre de la Libération}, notice «~Marcel Faure~». \url{https://www.ordredelaliberation.fr/fr/compagnons/marcel-faure}.
\item \textbf{PSS-Archi}, immeuble FR-31555-9853 (ancien bazar Labit, Toulouse). \url{https://www.pss-archi.eu/immeubles/FR-31555-9853.html}.
\item \textbf{RATP}, «~Histoire de la ligne~4~». \url{https://www.ratp.fr/decouvrir/patrimoine/histoire-metro-ligne-4}.
\item \textbf{Science Museum Group}, «~F~C Calvert and Company 1859--1965~». \url{https://collection.sciencemuseumgroup.org.uk/people/cp60291/f-c-calvert-and-company}.
\item \textbf{Spink}, \textit{Orders, Decorations and Medals e-Auction~3} (vente 21102), lot~188 : «~A Great War M.C. group of four awarded to Captain E.~M.~Ford, Royal Army Service Corps~». \url{https://spink.com/lot/21102000188}.
\item \textbf{The Canadian Encyclopedia}, «~Battle of Mount Sorrel~» et «~Battle of St. Eloi Craters~».
\item \textbf{The Long, Long Trail} (Chris \textsc{Baker}) [LLT], \url{https://www.longlongtrail.co.uk}, pages consultées directement ou par des copies datées de la Wayback Machine (\url{https://web.archive.org}) : «~18th (Eastern) Division~» (copie du 9~février 2017) ; «~31st Division~» ; «~1st Indian Cavalry Division (later 4th Cavalry Division)~» (copie du 8~mai 2023) ; «~2nd Indian Cavalry Division (later 5th Cavalry Division)~» ; «~The British Fifth or Reserve Army 1916--1918~» (copie du 23~juillet 2024) ; «~The British II Corps in France and Flanders~» (copie du 6~novembre 2024) ; «~The British IV Corps in France and Flanders~» (copie du 4~décembre 2021) ; «~The British V Corps in France and Flanders~» (copie du 5~août 2020) ; «~The British XIII Corps in France and Flanders~» ; «~Locations of British Casualty Clearing Stations~» (copie du 13~janvier 2018) ; «~The Army Service Corps in the First World War~», «~Army Service Corps Horse Transport Companies~» et «~Army Service Corps Mechanical Transport Companies~» (copie du 8~décembre 2021) ; «~The British Trench Mortar Batteries in the First World War~» ; «~Expeditionary Forces Canteens~» ; «~King's Liverpool Regiment~».
\item \textbf{Toulouse Métropole}, «~Balades nature~» (carte, 2022) et notices de quartier «~Rangueil, Sauzelong, Jules-Julien, Pech-David, Pouvourville~» et «~Maison des associations~». \url{https://metropole.toulouse.fr}.
\item \textbf{UNESCO}, Centre du patrimoine mondial, «~Beffrois de Belgique et de France~», n\textsuperscript{o}~943bis (Armentières 943-032, Bailleul 943-033, Lucheux 943-052). \url{https://whc.unesco.org/fr/list/943}.
\item \textbf{UNSW Canberra}, \textit{The AIF Project}, notices Henry Coates Wentworth (pid 318876), Fitzwilliam Wentworth (pid 318872, 318873), D'Arcy Wentworth (pid 318870) et autres Wentworth (pid 318878, 318879). \url{https://aif.adfa.edu.au}.
\item \textbf{Ville de Blois}, dossier de presse «~Les belles images de Poulain~», exposition du 15~septembre au 31~octobre 2015, 16~septembre 2015 (diffusé par l'APHG). \url{https://www.aphg.fr/IMG/pdf/dpexpositionpoulain.pdf}.
\item \textbf{Ville de Drancy}, exposition «~Drancy 14-18~», «~La gare régulatrice du Bourget~».
\item \textbf{Wikidata}, Q559663 (Louis-Fernand Flutre). \url{https://www.wikidata.org/wiki/Q559663}.
\item \textbf{Wikipédia} (points d'entrée, informations recoupées ailleurs) : en français, «~Ancien tramway de Toulouse~», «~Basilique Notre-Dame de Brebières~», «~Carnoy~», «~Chemins de fer départementaux de la Somme~» (§~«~Ligne d'Albert à Doullens~»), «~Cimetière de Terre-Cabade~», «~Corbeil-Essonnes~», «~Ernest Zyromski~», «~Estaires~», «~Gare de Berguette-Isbergues~», «~Ligne d'Arras à Dunkerque-Locale~», «~Ligne~1~», «~Ligne~4~» et «~Ligne~5 du métro de Paris~», «~Ligne de Villeneuve-Saint-Georges à Montargis~», «~Liste des monuments historiques de la Somme~», «~Mesnil-Martinsart~», «~Montaudran~», «~Permissionnaires français de la Première Guerre mondiale~», «~Pierre des Vallières~» ; en anglais, «~Bagne of Toulon~», «~Battle of Hallue~», «~Battle of Nagyszeben (1916)~», «~II ANZAC Corps~», «~Leipzig Salient~», «~Louis-Fernand Flutre~», «~Mines on the first day of the Somme~», «~Pygmalion (play)~», «~Réseau Albert~», «~Serre-lès-Puisieux~», «~Toulouse-Matabiau station~».
\end{biblist}

\endgroup
